\section{Introduction}\label{sec:intro}

The transition to mixed-autonomy traffic creates an organizational problem in addition to a control problem. Autonomous vehicles (AVs) will be dispatched by multiple firms, platforms, or routing authorities, while human-driven vehicles (HDVs) continue to respond to congestion through individual route choice. The firms need not value the same operational outcome: one may emphasize the travel time of its own fleet, another may internalize a larger share of congestion, and a third may encode reliability, energy, or service obligations. The same set of firms can therefore be operated as independent routing centers, coordinated in partial coalitions, or placed under a common dispatcher. These arrangements are not equivalent when the aggregate demand and the physical road network are held fixed.

This paper studies the resulting design problem. Given a set of heterogeneous
AV companies, when should their routing be centralized, when should partial
independence be retained, and can a local rule select the organization without
relying on a full ex post enumeration of partitions? The question is
deliberately conditional. A partition and a governance rule are specified
before routing begins; companies retain the OD obligations associated with the
trips they have accepted; and the analysis compares the equilibria induced by
alternative arrangements. We do not model whether firms merge or sign an
agreement. The object of interest is the operational consequence of an
arrangement that a platform, fleet manager, or public authority could choose
among feasible alternatives.

The need for this distinction is practical as well as theoretical. Ownership concentration statistics such as the Herfindahl--Hirschman Index (HHI) summarize who controls how much mass and are used in market-screening contexts \citep{usdojftc2023}. They do not record what those firms optimize, which firms are grouped together, how many routing decision centers remain, which OD obligations each member must preserve, or which parts of a coalition response can be offset by HDV rerouting. Consequently, two organizations with the same company HHI, coalition size, and aggregate demand can induce different network flows. The conditional design problem can be written as
\begin{equation}
\Pi^\star(\alpha,\theta,\gamma)\in\arg\min_{\Pi\in\mathfrak P(N)}J\!\left(x^\star(\Pi;\alpha,\theta,\gamma)\right),
\label{eq:design-map-intro}
\end{equation}
where $\Pi$ is a feasible partition of the companies, $\gamma$ is a governance intensity, $\alpha$ is AV penetration, $\theta$ collects objective types, and $x^\star$ is the mixed-autonomy equilibrium induced by that arrangement. The map is a design object, not an endogenous coalition-formation equilibrium.

We implement this design problem as a coalition-structured mixed Nash--Wardrop game. Each coalition is an atomic splittable player with a convex management objective, while HDVs form a nonatomic Wardrop population. All vehicles experience the same physical edge cost; heterogeneity enters through the objectives used to manage AV flows. The main treatment uses identity-preserving coordination: members may jointly choose routes, but each member must satisfy its own accepted OD demand. Operational pooling, which relaxes those member-level constraints, is analyzed separately because it changes feasible authority rather than merely changing objective aggregation. The formulation links atomic congestion games and Cournot--Wardrop interaction \citep{haurie1985,orda1993,harker1988,boulogne2002} with multiclass assignment \citep{dafermos1972} and mixed-autonomy routing \citep{mehr2019,lazar2021,biyik2021}.

The model makes the transmission mechanism explicit. Demand conservation
restricts every coalition's route response to an OD-feasible tangent. After
this response is aggregated in edge space, HDVs can absorb any component that
lies in their own Wardrop adjustment space. The remaining component is the
activated coalition response. Projecting that response onto the gradient of
system travel time produces a directional test for a proposed merge. The
criterion distinguishes a screened regime, in which organization is locally
neutral, from regimes in which stronger coupling raises or lowers congestion.
In the canonical bottleneck, it also identifies an interior effective-response
target at which partial coordination can dominate both grand and singleton
organization.

Our first contribution is a unified representation of heterogeneous coalition control in mixed-autonomy routing. The model separates the physical cost function shared by all vehicles from the company management objective, the coalition governance matrix, the number of underlying companies $n$, and the number of active coalition decision centers $K=|\Pi|$. The convex quadratic family contains independent routing, partial cross-member internalization, and common-dispatcher arrangements while preserving the OD obligations needed for a clean coalition comparison.

Our second contribution is a regime theory for coalition design. Starting from
the bordered sensitivity systems used for equilibrium networks with
polyhedral constraints \citep{tobin1988,qiu1989,patriksson2004,josefsson2006},
we derive the OD-restricted coalition response, aggregate it in edge space,
and quotient it by the endogenous HDV Wardrop tangent. On a fixed affine
support, the activated response signature is sufficient for aggregate
equilibrium and minimal for the complete conditional response map. Its
system-cost projection gives a general-network merge-direction score. The
canonical bottleneck sharpens this score into a threshold theorem: complete
HDV screening makes organization neutral, a boundary regime favors changes
that reduce effective response, and an interior regime can favor partial
coordination. Parametric equilibrium computation supplies the complementary
support-traversal perspective \citep{klimm2025}.

Our third contribution is a certified, regret-based selection framework. We
evaluate one- and two-step agglomerative policies against complete
four-company scans, using enumeration only to audit the selected partitions.
Both policies attain a state optimum in all four Sioux Falls conditions. On
independent Braess and directed-grid networks, the one-step rule is exact in
$29$ of $32$ states with mean regret $0.332$ recovery points, while two-step
lookahead is exact in all $32$ states. Across the $480$ certified
cross-network profiles, grand coordination is preferred at lower penetration,
whereas partial coordination emerges at high penetration under sufficiently
strong governance on both topologies. An explicit governance-burden model
produces calculable intervals in which grand, partial, or fragmented control is
optimal. Composition sensitivity, company-count, spatial-exposure, and HHI
diagnostics establish why the policy must use the response state rather than a
mass statistic.

The numerical evidence is deliberately bounded. The three networks are
computational benchmarks, not an empirical sample of mobility markets. The
reported rankings hold within the declared objective family,
identity-preserving governance rule, demand portfolios, and route sets; the
affine regime results are not asserted as global BPR theorems. The two-step
policy's zero observed regret is a benchmark result rather than a universal
guarantee. Pricing, passenger acceptance, empty-vehicle repositioning,
platform profits, empirical calibration, and endogenous formation remain
outside the paper's scope.

The remainder of the paper follows this logic. Section~\ref{sec:model} defines
the physical network, company objectives, coalition governance, equilibrium,
and performance measures. Section~\ref{sec:response} establishes
well-posedness, the activated response, the general merge score, and the
canonical regimes. Section~\ref{sec:ownership} interprets coalition
composition, governance burden, HHI, and the boundary between objective
coordination and pooling. Section~\ref{sec:experiments} reports mechanism,
regret, and cross-network evidence. Section~\ref{sec:discussion} discusses
implementation and scope, and Section~\ref{sec:conclusion} summarizes the
bounded implications.
